% Section 13 Header
\section{Pattern 11: Matrix Traversal (The "Islands" Pattern)}

\begin{frame}[c]{\empty}
    \begin{center}
        \Large\textbf{\secname}
    \end{center}
\end{frame}

% Slide 1: Concept
\begin{frame}{What is the Matrix Traversal Pattern?}
  \begin{itemize}
    \item \textbf{The Core Idea:} Treat a 2D grid (matrix) as an implicit graph where each cell $(r, c)$ is a node, and its adjacent neighbors (up, down, left, right) are connected by edges.
    \item \textbf{Traversal:} Use Depth First Search (DFS) or Breadth First Search (BFS) to explore connected regions.
  \end{itemize}
\end{frame}

\begin{frame}{Matrix Traversal Mechanics}
  \begin{itemize}
    \item \textbf{Boundary Check:} Always ensure row `r` and column `c` are valid:
      $$0 \le r < R \quad \text{and} \quad 0 \le c < C$$
    \item \textbf{Directions Array:} Use a utility array to iterate over neighbors cleanly:
      $$\text{dirs} = [(-1, 0), (1, 0), (0, -1), (0, 1)]$$
    \item \textbf{Visited Tracking:} Mark grid cells as visited (e.g., set `'1'` to `'0'` in-place if allowed, or use a `visited` boolean matrix) to prevent infinite loops.
  \end{itemize}
\end{frame}

% Problem 1: Number of Islands
\begin{frame}[fragile]{Problem: Counting Isolated Grid Clusters}
  \textbf{Scenario:}
  Given a 2D `matrix` of '1's (land) and '0's (water), count the number of isolated land clusters (islands) surrounded by water cardinally.
  
  \vspace{0.3cm}
  \textbf{Examples:}
  \begin{itemize}
    \item \textbf{Input:} \texttt{matrix = [['1','1','0'],['1','1','0'],['0','0','0']]} $\rightarrow$ \textbf{Output:} \texttt{1}
  \end{itemize}
\end{frame}

% Problem 2: Max Area of Island
\begin{frame}[fragile]{Problem: Largest Connected Grid Cluster}
  \textbf{Scenario:}
  Find the maximum area of a connected grid cluster (island) in a 2D `matrix`, where area is the number of cardinally connected land cells.
  
  \vspace{0.3cm}
  \textbf{Examples:}
  \begin{itemize}
    \item \textbf{Input:} \texttt{matrix = [[0, 1], [1, 1]]} $\rightarrow$ \textbf{Output:} \texttt{3}
  \end{itemize}
\end{frame}

% Problem 3: Rotting Oranges
\begin{frame}[fragile]{Problem: Spreading Contagion in Grid}
  \textbf{Scenario:}
  In a 2D `matrix` containing fresh items, contaminated items, and empty spaces, find the minimum time elapsed until all fresh items are contaminated. Spreads cardinally. Return -1 if impossible.
  
  \vspace{0.3cm}
  \textbf{Examples:}
  \begin{itemize}
    \item \textbf{Input:} \texttt{matrix = [[2,1,1],[1,1,0],[0,1,1]]} $\rightarrow$ \textbf{Output:} \texttt{4}
  \end{itemize}
\end{frame}

% Problem 4: Pacific Atlantic Water Flow
\begin{frame}[fragile]{Problem: Water Flow Boundaries}
  \textbf{Scenario:}
  Given a 2D `matrix` of heights, find all coordinates from which water can flow cardinally to both the left/top ocean (Pacific) and the right/bottom ocean (Atlantic).
  
  \vspace{0.3cm}
  \textbf{Examples:}
  \begin{itemize}
    \item \textbf{Input:} \texttt{matrix = [[1,2],[2,1]]} $\rightarrow$ \textbf{Output:} \texttt{[[0,0],[0,1],[1,0],[1,1]]}
  \end{itemize}
\end{frame}

% Common Triggers Slide
\begin{frame}{\secname\ -- Common Triggers}
  \begin{itemize}
    \item Counting distinct blocks (e.g. islands, provinces).
        \item Simulating cell expansion over time (e.g. rotting oranges, fire spreading).
        \item Finding paths or reachable regions from boundaries.
  \end{itemize}
\end{frame}
